% >>> [agent-codex 2026-11-27 compact-book]
\documentclass[11pt,oneside,openany]{amsbook}
\usepackage[T1]{fontenc}
\usepackage[utf8]{inputenc}
\usepackage{lmodern}
\usepackage{amsmath,amssymb,amsthm,mathtools}
\usepackage{booktabs,array,longtable}
\usepackage{xcolor}
\usepackage[hidelinks]{hyperref}
\emergencystretch=4em

\newtheorem{theorem}{Theorem}[section]
\newtheorem{proposition}[theorem]{Proposition}
\newtheorem{definition}[theorem]{Definition}
\newtheorem{question}[theorem]{Question}
\newtheorem{example}[theorem]{Example}
\theoremstyle{remark}
\newtheorem{remark}[theorem]{Remark}
\renewcommand{\thesection}{\thechapter.\arabic{section}}
\numberwithin{equation}{section}

\newcommand{\C}{\mathbb C}
\newcommand{\Q}{\mathbb Q}
\newcommand{\R}{\mathbb R}
\newcommand{\Z}{\mathbb Z}
\newcommand{\eps}{\varepsilon}
\newcommand{\Lam}{\Lambda}
\newcommand{\re}{\operatorname{Re}}
\newcommand{\covol}{\operatorname{covol}}
\newcommand{\Sha}{\operatorname{Sha}}
\newcommand{\todo}[1]{\textcolor{red}{[todo: #1]}}
\newcommand{\taylor}[1]{\textcolor{blue}{[Taylor: #1]}}
\newcommand{\agentreply}[1]{\textcolor{gray}{[agent: #1]}}
\newcommand{\ghroot}{https://github.com/lsernau/grh-lean/blob/}
\newcommand{\ghrev}{72fead102f663d5c7a407964a13dc0cf9d977207}
\newcommand{\leanfile}[2]{%
  \href{\ghroot\ghrev/grh-compact/src/#1}{\nolinkurl{#2}}}

\begin{document}
\begin{abstract}
This manuscript traces the compact Lean formalization of the
Riemann hypothesis and Dirichlet GRH to a pinned source commit. It
explains how bounded automorphic forms, nearly spanning families,
and integral height bounds yield a one-sided weighted prime-error
estimate; a positivity criterion then excludes zeros to the right
of the critical line. The elementary comparison is proved here.
The construction of the forms and the deep height estimates are
stated with source locators and remain obligations for a human
mathematical reading. The formal project has passed reported
machine checks; this AI-drafted exposition has not been refereed.
\end{abstract}
\title{The Compact Lean Argument for Dirichlet GRH}
\author{Prepared for Taylor Dupuy\\[2ex]
  Working draft, 27 November 2026\\[1ex]
  AI GENERATED --- MATHEMATICAL REVIEW IN PROGRESS}
\date{27 November 2026}
\maketitle
\tableofcontents

\chapter*{Introduction}
Levent Alp\"oge announced a formal proof of the Riemann hypothesis
and its Dirichlet $L$-function analogue on 20 November 2026
\cite[pp.~1--4]{Alpoge2026}. Luke Sernau, colleagues at Anthropic,
and Claude subsequently published a shorter Lean rewrite. Its compact
edition has 794 Lean source files and 76,149 nonblank, noncomment
lines in \texttt{src/}; the total is 112,827 lines. The readable
edition has 803 files and 299,902 lines of code. Both prove the same
five statements \cite[Summary]{Sutherland2026}.

The reader should distinguish the formal result from this account.
Sutherland reports completed kernel and independent checker runs for
both editions \cite[Summary and What was verified]{Sutherland2026}.
The authors of the compact edition state that no mathematician has
read the proof through and that no paper accompanies it
\cite[Status]{Sernau2026}. This manuscript explains the dependency
structure and checks the elementary comparison in its final argument.
It does not supply an independent proof of the deep inputs.

The order follows the earlier lecture notes \cite[Chs.~1--3]{Dupuy2026}:
the claimed theorem and an analogy, the entire argument assembled
from named inputs, then the constructions that supply those inputs.
The compact rewrite has different definitions and proof terms from
the original archive, so the existing book cannot simply be relabeled
as an explanation of this edition.

\subsection*{Software and AI assistance}
An OpenAI Codex desktop agent drafted this manuscript on
27 November 2026 \cite[\href{https://openai.com/codex}{product
page}]{OpenAICodex}. The runtime did not expose its exact model
identifier. The agent was given the earlier AI-generated book and
read the pinned compact Lean source and the original Alp\"oge note.
Its explanation follows those sources; the comparison in
Proposition~\ref{prop:comparison} unpacks the source's
\texttt{star\_inequality}, rather than claiming a new result.
The source search found no independent paper on the compact rewrite.
Responsibility for any released version rests exclusively with its
human author. No human has reviewed this draft's mathematical claims.

\chapter{The target and the route}
\label{ch:target}
\section{Exactly what the Lean statements say}
\label{sec:statements}
The source file \leanfile{Statements.lean}{Statements.lean} imports
only Mathlib. It fixes four propositions without reference to the
proof's internal definitions. The first final theorem has Mathlib's
type \texttt{RiemannHypothesis}; the other four have these types
\cite[\texttt{src/Statements.lean} and
\texttt{src/Main.lean}]{Sernau2026}.

Write $L(s,\chi)$ for Mathlib's Dirichlet character $L$-function.
The source quantifies over every positive modulus $N$ and every
complex valued Dirichlet character modulo $N$. It asserts
\begin{equation}
\label{eq:nonvanishing}
\re s>\tfrac12,\quad (\chi\ne1\text{ or }s\ne1)
\quad\Longrightarrow\quad L(s,\chi)\ne0.
\end{equation}
For $0<\re s<1$, it also asserts that a zero of $L(s,\chi)$ has
$\re s=\tfrac12$, both for arbitrary and for primitive characters.
For the Riemann zeta function it states that every zero other than
$-2,-4,-6,\ldots$ has real part $\tfrac12$; the formal predicate
also excludes $s=1$, the pole.

\begin{remark}
\label{rem:scope}
The Dirichlet statements do not discuss zeros with $\re s\le0$.
The primitive-character theorem is a specialization of the theorem
for all characters. The term ``GRH'' in this manuscript refers only
to Dirichlet $L$-functions.
\end{remark}

\section{The analytic reduction}
\label{sec:analytic-reduction}
The central declaration is \texttt{New.grh} in
\leanfile{New/Main.lean}{New/Main.lean}. It proves
\eqref{eq:nonvanishing}. The functional equation then reflects a
zero in $0<\re s<\tfrac12$ into $\re s>\tfrac12$; separate
zeta identities treat the nontrivial zeta zeros. This is the route
implemented in the final five declarations of
\leanfile{Main.lean}{Main.lean}. The reduction is logically useful:
the arithmetic construction only needs an upper bound on one side
of a weighted prime count. No two-sided prime number theorem in
progressions is assumed at the final step.

\chapter{A function-field comparison}
\label{ch:model}
\section{Why a bound on primes can control zeros}
The Stepanov--Bombieri argument for curves over finite fields counts
points by placing a high-order vanishing condition on an auxiliary
function. Its zero divisor and its pole divisor then constrain each
other. A contradiction arises if the count exceeds the square-root
error scale. The proof is developed in the earlier book
\cite[Ch.~2]{Dupuy2026}. The present chapter uses that argument only
as an analogy; it is not an input to the compact Lean derivation.

The compact proof replaces a single auxiliary function by a family
of automorphic forms and their theta lifts. The forms have an upper
size bound. A global integrality and product-formula argument gives
a lower bound for the same covolume. The excess between the bounds
is a weighted prime sum. The comparison resembles a degree count
in the function-field argument, although the objects and the
estimates are substantially different.

\begin{remark}
\label{rem:analogy-limit}
The estimate for a curve over $\mathbb F_q$ does not imply the
arithmetic estimate used below. In the compact development the
automorphic, Galois, $p$-adic, and height inputs are proved in Lean
files under \texttt{src/New} and \texttt{src/Lib}; this chapter
supplies none of those arguments.
The earlier book shows that the number-field one-sided prime
bound corresponds to a lower point-count bound, while the
Stepanov--Bombieri argument gives an upper point-count bound
\cite[Ch.~2, direction of the bound]{Dupuy2026}.
Its alternative ``sign-only'' argument using $qN=M^k+R$
is conditional on a hypothesis that fails for some values of
$k$ in the simplest quadric cases
\cite[Ch.~2, second hypothesis]{Dupuy2026}.
Neither analogy supplies a missing step in Theorems A--C.
\end{remark}

\chapter{The argument assembled}
\label{ch:assembly}
\section{The data and the two bounds}
Fix admissible parameters $P$ and data $D$ as in
\leanfile{New/Data/DataExist.lean}{DataExist.lean}. The main scale is
$x$, the later scale is $x_1$, the large auxiliary parameter is $b$,
and $m$ is the number of theta lifts. Put
\begin{equation}
\label{eq:budget}
f=\varphi(N_0)\,\#\mathrm{tops},\qquad
\mathcal B=f a(\log b)^3,\qquad
L=(\log b)\mathcal B.
\end{equation}
Here $a$ is the long order, $\mathrm{tops}$ is a set of infinite
places, and $D_{\rm row}$ is the dimension of the row space.
Write $h_0=m k\mathcal U$ for the main height term and
$R$ for the remainder of the prime sum after its expected terms
are removed. Chapter~\ref{ch:prime-error} gives $R$ explicitly.

Theorems A and B of the compact source are respectively
\texttt{Data.short\_forms} and \texttt{Data.enough\_forms} in
\leanfile{New/Endgame/ShortAndEnough.lean}{ShortAndEnough.lean}.
They give independent rows $y_1,\ldots,y_n$ with
\begin{align}
\label{eq:nearly-all}
n&\ge(1-4/b^2)D_{\rm row},\\
\label{eq:index}
\delta&\le5D_{\rm row}\log(1+b^2)\mathcal B,
\end{align}
where $\delta$ controls the index of the integral lifts.
For each complex embedding $\sigma$, their squared covolume obeys
\begin{equation}
\label{eq:upper-covolume}
\covol_\sigma^2(\Theta_1 y_1,\ldots,\Theta_1 y_n)
\le\exp\bigl(n(h_0-mR+L)\bigr).
\end{equation}
This exact combined output is
\texttt{Data.exists\_family\_covolumeSq\_le} in
\leanfile{New/Endgame/LongForms.lean}{LongForms.lean}.

Theorem C is \texttt{Data.integral\_forms\_long} in the same file.
For any independent nearly spanning family satisfying the index
condition and a squared-covolume upper bound $e^V$, it gives
\begin{equation}
\label{eq:lower-covolume}
D_{\rm row}(h_0-3L)\le V+2\delta.
\end{equation}
This is a lower bound for the logarithmic covolume, despite the
declaration's name. Its proof uses the heights of twisted volume
ratios and the nonnegativity of the height of $1$.

\section{The cancellation}
\label{sec:cancellation}
The formal comparison appears as \texttt{star\_inequality} in
\leanfile{New/Main.lean}{New/Main.lean}. The following calculation
records its mathematical content. The source also proves the
crude estimate $|h_0|+m|R|\le b^2x^{99}$.

\begin{proposition}
\label{prop:comparison}
Let $D$ be admissible data satisfying
\eqref{eq:nearly-all}--\eqref{eq:upper-covolume}, the lower bound
\eqref{eq:lower-covolume}, and
$|h_0|+m|R|\le b^2x^{99}$. Then
\begin{equation}
\label{eq:comparison-raw}
mR\le4L+4x^{99}
      +10\log(1+b^2)\mathcal B.
\end{equation}
With the admissible scale inequalities used in
\texttt{star\_inequality}, the right side is at most
$(\log b)^2\mathcal B$.
\end{proposition}
\begin{proof}
Set $D_0=D_{\rm row}$ and $h=h_0-mR+L$.
The upper bound gives $V=nh$. Apply \eqref{eq:lower-covolume}:
\[
D_0(h_0-3L)\le nh+2\delta
 =D_0h-(D_0-n)h+2\delta.
\]
Since $D_0-n\le4D_0/b^2$ and $-h\le b^2x^{99}$, division by
$D_0>0$ gives
\[
h_0-3L\le h_0-mR+L+4x^{99}+2\delta/D_0.
\]
The index bound \eqref{eq:index} yields
\eqref{eq:comparison-raw}. The final scale comparison is
the monomial bound \texttt{D.le\_of\_hasLogBound} in the source.
\end{proof}

\begin{remark}
\label{rem:comparison-status}
The proof of Proposition~\ref{prop:comparison} checks the algebra
between A, B, and C. It does not establish those three theorems,
the existence of the admissible data, or the numerical scale bounds.
Their Lean proofs are separate dependencies.
\end{remark}

\section{From the comparison to a one-sided error bound}
\label{sec:one-sided}
The choice of $D$ includes
\begin{equation}
\label{eq:remainder-identity}
R=a\bigl(T(R_{N_0}(x)-R_{N_0}(x_1))
               +H R_1(x_1)\bigr),
\end{equation}
where $T=\#\mathrm{thick}$, $H=\#\mathrm{tops}$, and
$R_M$ is the normalized weighted prime error defined in
Section~\ref{sec:error-definition}. This is the exact identity
\texttt{Data.remainder\_eq}.

The window theorem in Section~\ref{sec:window} permits $x_1$ to be
chosen with
\[
R_{N_0}(x_1)-(H/T)R_1(x_1)\le x^{-10^6}.
\]
Divide \eqref{eq:remainder-identity} by $maT$ and use
Proposition~\ref{prop:comparison}. With $m\ge\sqrt{x}$,
$T\ge H/\log b$, and $\log b=x^\eps$, the formal scale estimates
give, for $0<\eps<1/100$,
\begin{equation}
\label{eq:error-bound}
R_{N_0}(x)\le x^{8\eps-1/2}
\qquad\text{for all sufficiently large }x.
\end{equation}
\begin{proposition}
\label{prop:one-sided-scale}
Let $0<\eps<1/100$, and suppose $D$ satisfies
\eqref{eq:remainder-identity}, the star inequality
$mR\le f a(\log b)^5$, and
$m\ge\sqrt x$, $T\ge H/\log b$, $f=H\varphi(N_0)$,
$\log b=x^\eps$. Suppose also that its second scale
satisfies
$R_{N_0}(x_1)-(H/T)R_1(x_1)\le x^{-10^6}$.
Then
\[
R_{N_0}(x)\le
 \varphi(N_0)x^{6\eps-1/2}+x^{-10^6}.
\]
For fixed $N_0$, this implies \eqref{eq:error-bound}
when $x$ is sufficiently large.
\end{proposition}
\begin{proof}
Divide \eqref{eq:remainder-identity} by $maT$ and
rearrange. The star inequality and the second-scale
bound give
\[
R_{N_0}(x)\le\frac{f(\log b)^5}{mT}
 +x^{-10^6}.
\]
Since $f/T\le\varphi(N_0)\log b$, the first term is
at most $\varphi(N_0)(\log b)^6/\sqrt x$.
For $\eps>0$, the fixed factor $\varphi(N_0)$ and
the smaller power $x^{-10^6}$ are absorbed by
$x^{2\eps}$ for large $x$.
\end{proof}
The exponent and the quantifiers are those of
\path{eventually_primeError_le_rpow} in
\leanfile{New/Main.lean}{New/Main.lean}.

\section{The place count behind the remainder identity}
\label{sec:remainder-proof}
The identity \eqref{eq:remainder-identity} is a finite
count, followed by cancellation of its explicit main
terms. It is useful to show that count, since it is where
the modulus $N_0$ first enters the prime sum. Let
$H=\#\mathrm{tops}$ and $T=\#\mathrm{thick}$.
For a positive integer $n$, set
\[
C(n)=\#\{\tau\in\mathrm{tops}:
       (-n)\tau\in\mathrm{thick}\}.
\]
The top places form the class $1\pmod{N_0}$ and the
thick places form the class $-1\pmod{N_0}$.
Multiplication by $-n$ therefore yields, after
multiplication by the weight $\lambda_{d,1}(n)$,
\begin{equation}
\label{eq:place-count}
C(n)\lambda_{d,1}(n)
 =T\lambda_{d,N_0}(n).
\end{equation}
If $n\equiv1\pmod{N_0}$, this multiplication permutes
the tops onto the class containing the thick places;
exactly $T$ top places map into them. Otherwise no top
place maps into the thick class. The source handles the
nonunit case through its \texttt{placeOfInt} definition.
Equation~\eqref{eq:place-count} is the local counting
claim \texttt{A} inside \texttt{Data.remainder\_eq} in
\leanfile{New/Endgame/ShortAndEnough.lean}%
{ShortAndEnough.lean}.

At a top place $\tau$, the nominal weight of
$(-n)\tau$ is $a/x$ when it is thick and $a/x_1$
otherwise. Hence the sum of the terms in
\eqref{eq:data-prime-sum} over the tops equals
\begin{align}
\label{eq:asks-sum}
\sum_{\tau\in\mathrm{tops}}A(a,\tau)
={}&T\sum_{n\le a}\lambda_{d,N_0}(n)
  \left[\left(\frac an-\frac ax\right)^+
        -\left(\frac an-\frac a{x_1}\right)^+\right]\\
 &+H\sum_{n\le a}\lambda_{d,1}(n)
       \left(\frac an-\frac a{x_1}\right)^+.
\end{align}
Apply \eqref{eq:finite-error} at $x$ and $x_1$.
The constants $\re\mathcal L_{d,N_0}(1)$ cancel
between the first two sums. By
\eqref{eq:mean-at-one}, the last sum has constant
$-1-\operatorname{CS}(d)$. Thus
\begin{align}
\label{eq:asks-expanded}
\sum_{\tau\in\mathrm{tops}}A(a,\tau)
=a\bigg[&T\bigl(R_{N_0}(x)-R_{N_0}(x_1)\bigr)
       +H R_1(x_1)\\
 &+\frac{T}{\varphi(N_0)}(\log x-\log x_1)
       +H(\log x_1-1-\operatorname{CS}(d))\bigg].
\end{align}

There are $f=H\varphi(N_0)$ infinite places, of which
$T$ are thick. Consequently the mean logarithmic
nominal weight is
\[
\bar w=\log a-
 \frac{T\log x+(f-T)\log x_1}{f}.
\]
Substituting this in \eqref{eq:main-ask} gives
\begin{equation}
\label{eq:ask-cancellation}
H A_{\rm main}=a\left[
 \frac{T}{\varphi(N_0)}(\log x-\log x_1)
 +H(\log x_1-1-\operatorname{CS}(d))\right].
\end{equation}
By definition $R=\sum_{\tau\in\mathrm{tops}}
(A(a,\tau)-A_{\rm main})$. Subtracting
\eqref{eq:ask-cancellation} from
\eqref{eq:asks-expanded} proves
\eqref{eq:remainder-identity}. This derivation uses
the source's exact finite count and the already proved
finite-error identity; no asymptotic prime theorem is
hidden in the cancellation.

\chapter{Prime errors and the analytic endgame}
\label{ch:prime-error}
\section{The weighted count}
\label{sec:error-definition}
Let $d$ be a prime congruent to $7$ modulo $8$ and coprime to $M$.
For the associated quadratic character $\chi_d$, the compact
source defines a nonnegative weight on positive integers,
\begin{equation}
\label{eq:lambda}
\lambda_{d,M}(n)=
  \mathbf1_{n\equiv1\,({\rm mod}\,M)}
  \Lam(n)(1-\chi_d(n)).
\end{equation}
The Legendre factor vanishes at split primes and doubles the weight
at nonsplit primes. The von Mangoldt factor also includes prime
powers. In particular, ``prime sum'' is a conventional shorthand
for a weighted prime-power sum.

Let
\[
S_r(y)=y^{-r}\sum_{1\le n\le y}\lambda_{d,M}(n)n^{r-1}
\quad(r=0,1).
\]
The normalized error is
\begin{equation}
\label{eq:prime-error}
R_M(y)=S_0(y)-S_1(y)-\frac{\log y}{\varphi(M)}
       -\re\mathcal L_{d,M}(1),
\end{equation}
where $\mathcal L_{d,M}$ is a mean of differences of logarithmic
derivatives of pole-removed Dirichlet $L$-functions. The precise mean
is \texttt{logDerivMean} in
\leanfile{New/Endgame/PrimeError.lean}{PrimeError.lean}.
For $0<y\le A$, the source proves the finite identity
\begin{equation}
\label{eq:finite-error}
R_M(y)=\frac1A\sum_{1\le n\le A}\lambda_{d,M}(n)
          (A/n-A/y)^+
      -\frac{\log y}{\varphi(M)}
      -\re\mathcal L_{d,M}(1),
\end{equation}
where $u^+=\max(u,0)$.
Equation~\eqref{eq:finite-error} is the interface between the
divisibility estimates for forms and the analytic prime error.

\begin{proposition}
\label{prop:finite-error}
Let $d$ be a prime, let $M\ge1$, and let $0<y\le A$ with
$A$ a positive integer. For $\lambda_{d,M}$ and $R_M$ defined
by \eqref{eq:lambda} and \eqref{eq:prime-error}, the finite
identity \eqref{eq:finite-error} holds.
\end{proposition}
\begin{proof}
The definitions give
\[
S_0(y)-S_1(y)
=\sum_{1\le n\le y}\lambda_{d,M}(n)
   \left(\frac1n-\frac1y\right).
\]
For $n\le y$, the expression in parentheses equals
$(A/n-A/y)^+/A$. For $y<n\le A$ the positive part is zero.
Appending those zero terms gives \eqref{eq:finite-error}.
\end{proof}

\section{A good point in a long window}
\label{sec:window}
The source's predicate \texttt{IsPolar} describes the analytic
continuation and pole structure of a Laplace transform. The file
\leanfile{New/Endgame/OneSidedCriterion.lean}%
{OneSidedCriterion.lean} proves a window statement in this setting.

\begin{definition}
\label{def:is-polar}
Let $E:\R\to\R$ be measurable and let $r:\C\to\C$.
The source predicate $\operatorname{IsPolar}(E,r)$ means that
$|E(v)|\le C e^v$ for $v\ge0$ and some $C$, and that
the Laplace transform
\[
\Phi(s)=\int_0^\infty E(v)e^{(1-s)v}\,dv
\quad(\re s>2)
\]
extends meromorphically. The function $(s-1)\Phi(s)$
must be analytic on $\re s\ge1$. At every $\rho$ with
$\re\rho>0$, the limit of $(s-\rho)\Phi(s)$ as
$s\to\rho$ through $s\ne\rho$ must equal $r(\rho)$.
\end{definition}
This is \texttt{IsPolar} in
\leanfile{Lib/Analytic/Landau/PolarLandau.lean}%
{PolarLandau.lean}.

\begin{theorem}
\label{thm:window}
Let $\operatorname{IsPolar}(E,r)$ and
$\operatorname{IsPolar}(E',r')$ hold. For all sufficiently
large $t$ and every
$0\le X\le e^t$, there is $u\in[t,t^2]$ such that
\begin{equation}
\label{eq:window}
E(u)-XE'(u)\le e^{-t}.
\end{equation}
\end{theorem}
The statement is
\texttt{IsPolar.exists\_window\_sub\_mul\_le}. The source
proves it with a high moment of a Laplace transform and an
inequality that would be contradicted if \eqref{eq:window} failed
throughout the window. Here $E(u)=R_{N_0}(e^u)$ and
$E'(u)=R_1(e^u)$. The polar hypotheses are established for these
functions in \texttt{isPolar\_primeError}. A full paper proof of
the moment argument is still needed for this manuscript.

The quantitative step is more informative than the phrase
``sign change.'' For $R>3$ sufficiently near $3$, the source
obtains, for all large $t$ and every integer $k\ge0$,
\begin{equation}
\label{eq:moment}
\left|\int_0^\infty v^k e^{-3v}E(v)\,dv\right|
   \le\frac{k!\,t}{R^k},
\qquad |E(v)|\le t e^v\quad(v\ge0).
\end{equation}
The same holds for $E'$. Its proof differentiates the polar
Laplace transform under the integral and bounds the derivatives
by analyticity in a disk around $3$. To prove the window theorem,
the Lean source chooses $k=\lceil t^2/9\rceil$. If
$E(u)-XE'(u)>e^{-t}$ throughout $[t,t^2]$, positivity of a
weighted integral over $v>0$ gives a lower bound. The moments
\eqref{eq:moment} and two explicit exponential correction
terms give a strictly smaller upper bound. Thus the point
$u$ in \eqref{eq:window} exists. The constants $3$, $9$,
and the window $[t,t^2]$ belong to the formal statement and
its proof; they are not a general sign-change principle.

\section{The Landau criterion}
\label{sec:landau}
The formal lemma
\texttt{entireL\_ne\_zero\_of\_primeError\_le}
uses a one-sided bound such as \eqref{eq:error-bound}. It requires
the bound for every prime $d$ in a unit residue class that also
forces $d\equiv7\pmod8$. If $\beta>1/2$ and eventually
$R_M(x)\le x^{\beta-1}$ for those $d$, then no pole-removed
$L$-function for a character modulo $M$ has a zero in
$\re s>\beta$. Its proof combines a Landau positivity argument
with a quadratic twist chosen to avoid a possible zero. The
quadratic-twist input lives in
\leanfile{Lib/Analytic/QuadraticTwist/ZeroFreeTwist.lean}%
{ZeroFreeTwist.lean}. We do not replace that input by a blanket
claim that a one-sided bound for one chosen $d$ suffices.

For a character modulo $q$, \texttt{New.grh} passes to
$N_0=18480q$ and the unit class $N_0/2-1$. It applies
\eqref{eq:error-bound} for each prime $d$ in that class, takes
$\beta=1/2+8\eps$ below the real part of a proposed zero,
and invokes the criterion. This proves the half-plane
nonvanishing statement \eqref{eq:nonvanishing} within the Lean
development.

\section{Logarithmic derivatives and residues}
\label{sec:log-derivative}
The constant removed in \eqref{eq:prime-error} has an exact
analytic definition. Let $\widetilde L(s,\xi)$ be
\texttt{entireL}: it equals $L(s,\xi)$ for a nonprincipal
character and continues $(s-1)L(s,\xi)$ for a principal
one. This is the definition in
\leanfile{Lib/General/NumberTheory/QuadField.lean}%
{QuadField.lean}. With $\chi_d(n)=(n/d)$, the source sets
\begin{equation}
\label{eq:log-derivative-mean}
\mathcal L_{d,M}(s)=\frac1{\varphi(M)}
\left(\sum_{\xi\bmod M}
\left\{\frac{\widetilde L'}{\widetilde L}
  (s,\xi\chi_d)-
  \frac{\widetilde L'}{\widetilde L}(s,\xi)\right\}-1\right).
\end{equation}
The final $-1$ lies outside the character sum. At modulus
one, \texttt{re\_logDerivMean\_one} gives
\begin{equation}
\label{eq:mean-at-one}
\re\mathcal L_{d,1}(1)
  =-1-\operatorname{CS}(d),
\end{equation}
where $\operatorname{CS}(d)$ is the source's
Chowla--Selberg constant. Averaging the Dirichlet series
for a logarithmic derivative over $\xi\bmod M$ imposes
$n\equiv1\pmod M$. Subtracting the quadratic twist
produces $1-\chi_d(n)$. The principal pole contributes the
$\log y/\varphi(M)$ term. The precise identities are
\texttt{logDerivMean} and \texttt{primeError\_eq\_sum}
in \leanfile{New/Endgame/PrimeError.lean}{PrimeError.lean}.

The class \texttt{IsPolar} requires a real measurable
$E(v)=O(e^v)$ and a meromorphic continuation $\Phi$ of
\begin{equation}
\label{eq:laplace-transform}
\Phi(s)=\int_0^\infty E(v)e^{(1-s)v}\,dv
\qquad(\re s>2).
\end{equation}
It requires $(s-1)\Phi(s)$ analytic on $\re s\ge1$ and
specifies the residue at each $s$ with $\re s>0$.
See \leanfile{Lib/Analytic/Landau/PolarLandau.lean}%
{PolarLandau.lean}. For $E(v)=R_M(e^v)$, the proof of
\texttt{isPolar\_primeError} identifies the residue as
\begin{equation}
\label{eq:polar-residue}
r_{d,M}(s)=\frac{1}{\varphi(M)s(1-s)}
\sum_{\xi\bmod M}
\bigl(\operatorname{ord}_s\widetilde L(\xi)
  -\operatorname{ord}_s\widetilde L(\xi\chi_d)\bigr).
\end{equation}
Here $\operatorname{ord}_s$ is zero at a nonzero value and
the analytic order at a zero. The hypotheses are
$d\equiv7\pmod8$ and $(d,M)=1$. The local identity
\[
\frac{\widetilde L'}{\widetilde L}(z,\xi)
 =\frac{\operatorname{ord}_s\widetilde L(\xi)}{z-s}
 +\text{a holomorphic term}
\]
explains the difference of orders. The formal proof also
checks growth and removal of the pole at $s=1$; those
facts do not follow from the displayed local identity.

\section{Why moments force a window point}
\label{sec:moment-argument}
The moment estimate \eqref{eq:moment} follows by
differentiating \eqref{eq:laplace-transform} under the
integral. On $\re s>2$, the growth condition dominates
each differentiated integrand, giving
\begin{equation}
\label{eq:moment-derivative}
\Phi^{(k)}(s)=\int_0^\infty
 (-v)^kE(v)e^{(1-s)v}\,dv.
\end{equation}
The singularity at $1$ is removable under \texttt{IsPolar}.
One may therefore choose a disk centered at $4$ of radius
$R>3$ and use Cauchy estimates on that disk. For large
$t$ they give $|\Phi^{(k)}(4)|\le k!t/R^k$ for every $k$,
which is \eqref{eq:moment} by
\eqref{eq:moment-derivative}. The exact differentiation,
disk, and Taylor-series argument is
\texttt{IsPolar.eventually\_moment\_le} in
\leanfile{New/Endgame/OneSidedCriterion.lean}%
{OneSidedCriterion.lean}.

To see the window mechanism, suppose instead that
$E(u)-XE'(u)>e^{-t}$ at every $u\in[t,t^2]$.
Choose $k=\lceil t^2/9\rceil$. The formal proof makes
the following integrand nonnegative for $v>0$:
\begin{equation}
\label{eq:window-integrand}
v^k\left(e^{-Rv+7t}+e^{-v+2t-t^2}
  +e^{-3v}(E(v)-XE'(v)-e^{-t})\right).
\end{equation}
On the window, the last term is positive. Outside it,
the first two exponentials dominate its possible negative
part, using $|E(v)|,|E'(v)|\le te^v$ and $X\le e^t$.
Integration would give a nonnegative result. The moment
bound and
\begin{equation}
\label{eq:laplace-power}
\int_0^\infty v^k e^{-av}\,dv=\frac{k!}{a^{k+1}}
\qquad(a>0)
\end{equation}
give a negative upper bound for large $t$, a contradiction.
This describes the proof of
\texttt{IsPolar.exists\_window\_sub\_mul\_le}; its
numerical comparison of the three terms is a formal-only
input here.

\section{The one-sided Landau step}
\label{sec:landau-detail}
The general lemma \texttt{IsPolar.landau} says that an
eventual upper bound $E(v)\le e^{(\beta-1)v}$, together
with nonnegative real residues at real points $u>\beta$,
forces every residue with $\re s>\beta$ to vanish.
Its proof subtracts $E$ from a larger positive exponential
and applies the boundary-of-convergence argument for the
Laplace transform of a nonnegative density. This is the
Landau--Pringsheim mechanism in
\leanfile{Lib/Analytic/Landau/PolarLandau.lean}%
{PolarLandau.lean}.

The difference of zero orders in \eqref{eq:polar-residue}
does not have a sign by itself. The lemma
\texttt{exists\_entireL\_quadTwist\_ne\_zero} in
\leanfile{Lib/Analytic/QuadraticTwist/ZeroFreeTwist.lean}%
{ZeroFreeTwist.lean} has the following stronger output.
\begin{theorem}[Zero-free quadratic twist in the source]
\label{thm:zero-free-twist}
Let $\beta>1/2$, let $Y\ge0$, and let $a$ be a unit
residue class modulo $M$ such that every prime $d$ in
that class is congruent to $7\pmod8$ and coprime to
$M$. There is a prime $d\equiv a\pmod M$ for which
\[
\widetilde L(s,\xi\chi_d)\ne0
\quad\text{for every }\xi\bmod M,
\quad\re s\ge\beta,
\quad|\operatorname{Im}s|\le Y.
\]
\end{theorem}
The source proves this through a zero-density estimate
for quadratic twists averaged over the prime $d$.
If every prime in the class were bad, the resulting
bound would make the reciprocal-prime series in that
class summable. The source contradicts this with the
divergence of that series. Its zero-density theorem and
the arithmetic of twist conductors are substantial
formal inputs; the preceding argument identifies the
contradiction but does not reproduce those estimates.

Take $Y$ at least the imaginary part of a proposed
zero $\rho$. Then Theorem~\ref{thm:zero-free-twist}
selects a prime whose twisted factors avoid it and
all real points to the right of $\beta$.
The one-sided hypothesis holds for \emph{every} such prime,
so it holds for that $d$. The residue then detects the
untwisted zero. This explains the universal quantifier in
Section~\ref{sec:landau}.

\section{The final modulus and the critical line}
\label{sec:final-modulus}
To apply the one-sided criterion to a character modulo
$q$, the source chooses $N_0=18480q$ and the residue
class $a=N_0/2-1=9240q-1$ modulo $N_0$.
Since $4\mid N_0$,
\[
a^2-1=N_0(N_0/4-1),
\]
so $a$ is a unit. Since $8\mid9240$, every prime
$d\equiv a\pmod{N_0}$ satisfies $d\equiv7\pmod8$.
It is also coprime to $N_0$. These two elementary
facts supply the congruence hypotheses for the
quadratic-twist criterion.

Suppose $L(\rho,\chi)=0$ with $\re\rho>1/2$, away
from the principal pole. Put
\[
\eps=\frac{\min(\re\rho,1)-1/2}{51},
\qquad\beta=1/2+8\eps.
\]
Then $0<\eps<1/100$ and $\beta<\re\rho$.
Changing the level of $\chi$ from $q$ to $N_0$
preserves the relevant zero of the pole-removed
$L$-function, as checked in \texttt{New.grh} by
Mathlib's level-change identities. For every prime
$d$ in the class $a$, Proposition~\ref{prop:one-sided-scale}
and \texttt{exists\_data} supply the eventual bound
$R_{N_0}(x)\le x^{\beta-1}$. The criterion of
Section~\ref{sec:landau} contradicts the zero at
$\rho$. This is the proof of \texttt{New.grh} in
\leanfile{New/Main.lean}{Main.lean} once the deep
inputs have been supplied.

The final declaration \path{New.zeros_on_line}
uses a functional equation to reflect Dirichlet zeros.
The declaration \path{New.zeta_zeros_on_line} does the
same for zeta zeros left of $1/2$, moving them into the
excluded right half-plane. They also check the gamma
factor and pole or trivial-zero exceptions. Thus the
half-plane statement leads to the exact five target
types in \leanfile{Main.lean}{Main.lean}. The English
argument here records that reduction; its gamma-factor
identities remain imported from the formal source.

\chapter{The admissible data}
\label{ch:data}
\section{Scales and fields}
The finite data are arranged to make the comparison in
Chapter~\ref{ch:assembly} effective. The guide specifies
$18480\mid N_0$, a prime $d\equiv N_0/2-1\pmod{N_0}$,
$K=\Q(\sqrt{-d})$, and a cyclotomic extension $E=K(\zeta)$.
The ranks $k,r$ obey $r+1=k+m$, with
\begin{equation}
\label{eq:scales}
\sqrt{x}\le m\le2\sqrt{x},\quad
x^5\le r\le3x^5,\quad
(\log b)^{100}\le x,\quad
x^{10^6}\le x_1,\quad x_1^{10^6}\le b.
\end{equation}
The number of top infinite places is around $x^{11}$ and the
thick subset has at least $\#\mathrm{tops}/\log b$ elements.
There are also two primes whose product $g$ lies between
$b^{20}$ and $b^{22}$, and a level prime $N$ with
$\log N$ in a prescribed short window. These are features of
the data predicate, not independent choices one can make later.

The existence theorem \texttt{exists\_data} in
\leanfile{New/Data/DataExist.lean}{DataExist.lean} gives such
data for all sufficiently large $x$ when $0<\eps<1/100$.
In particular it fixes $\log x_1=x^{\eps/6}$ and
$\log b=x^\eps$. Its proof uses a Behrend-type set for
several equations and a pigeonhole argument in the class group.
The other ingredients include primes in progression and
construction of the field and place data. The source's precise
quantifiers are essential to the window choice in
Section~\ref{sec:one-sided}.

\section{The full scale record and its weights}
\label{sec:data-record}
The record \texttt{Data} in
\leanfile{New/Data/AdmissibleData.lean}{AdmissibleData.lean}
stores the logarithms of $x,x_1,b$, not merely their
values. Besides \eqref{eq:scales}, it requires
\begin{align}
\label{eq:log-scale-record}
10^{12}c&\le\log\log b,
&100\log\log b&\le\log x,\\
10^6\log x&\le\log x_1,
&10^6\log x_1&\le\log b,
\end{align}
and $6\log\log x_1\le\log\log b$. The constant $c$
absorbs quantities fixed by $N_0,d$, including a bound
from the cyclotomic generator construction. These
inequalities make every later comparison an explicit
comparison of logarithmic exponents.

The exponent $s$ in the conductor $N_1$ is linked to the
power-of-two rank $m$; the rank $r$ is a power of three.
The auxiliary field is $E=K(\zeta_{N_1})$, with
$N_1=3^c2^sN_0$ in the source's notation. The rank
$k=r+1-m$, the step size $t=5r$, and the ranks of the
intermediate spaces are
\begin{equation}
\label{eq:chain-ranks}
n_j=k+j(t+1),\qquad 0\le j\le m.
\end{equation}
The source also uses the gap $h=2r^8$ between successive
weights. These formulas are the declarations
\texttt{Data.k}, \texttt{Data.t}, \texttt{Data.n}, and
\texttt{Data.h} in \texttt{AdmissibleData.lean}.

The long order $a$ is the least multiple of $N_1$ at least
$xb$. The initial weights are the least multiples of $N_1$
at least $a/x$ and $a/x_1$. For an infinite place $\sigma$,
the weight at stage $j$ is
\begin{equation}
\label{eq:weight-stages}
w_j(\sigma)=
\begin{cases}
 b_1+jh+r^8,&\sigma\in\mathrm{thick},\\
 b_2+jh+r^8,&\sigma\notin\mathrm{thick}.
\end{cases}
\end{equation}
The nominal, unrounded weight is $a/x$ on thick places
and $a/x_1$ elsewhere. The rounding keeps integrality
conditions compatible with the conductor; the large gap
separates weights at successive lifts. The source's
\texttt{Data.primeSum} is the exact finite expression
\begin{equation}
\label{eq:data-prime-sum}
A(y,\sigma)=\sum_{1\le n\le a}
 \Lambda(n)(1-\chi_d(n))
 \left(\frac yn-
 w_{\rm nom}((-n)\sigma)\right)^+.
\end{equation}
The cost of one direction is
\begin{equation}
\label{eq:direction-cost}
\mathcal U=\sum_{\sigma\in\Phi}
\left(\log(w_0(\sigma)!)-
 w_0(\sigma)\log(\Omega/g)-
 A(w_{\rm nom}(\sigma),\sigma)\right),
\end{equation}
where $\Omega$ is \texttt{Parameters.bigOmega} and
$g=g_0g_1$. Equations \eqref{eq:data-prime-sum}
and \eqref{eq:direction-cost} port \texttt{Data.primeSum},
\texttt{Data.jetSize}, and \texttt{Data.dirCost}. They show
where the prime sum enters the height before the
covolume comparison.

\section{A large set avoiding several equations}
\label{sec:behrend}
The thick places must have two properties at once: they
occupy a substantial fraction of the tops, and they avoid
relations that would make different terms of the theta
chain collide. The combinatorial input is
\texttt{exists\_three\_term\_free} in
\leanfile{New/Data/DataExist.lean}{DataExist.lean}.

\begin{theorem}[Behrend set in the compact source]
\label{thm:behrend}
Let $\varphi:G\to H$ be a homomorphism of commutative
groups, and let $c,g\in G$. Let $n,M$ be positive
integers and $L\ge2$ with
$(LM)^n\le\operatorname{ord}(\varphi(g))$.
There is a finite set $T\subset c\langle g\rangle$ with
$|T|\ge M^n/(nM^2)$ such that for every
$0<i<j\le L$ and $\sigma,\tau,\rho\in T$,
\[
\varphi(\tau/\sigma)^j
 =\varphi(\rho/\sigma)^i
 \quad\Longrightarrow\quad
\varphi(\tau/\sigma)=1.
\]
\end{theorem}
\begin{proof}
Put $B=LM$. Among the $M^n$ vectors of
$\{0,\ldots,M-1\}^n$, choose a sphere $S$ of fixed
squared Euclidean norm containing at least
$M^n/(nM^2)$ vectors: the squared norm has fewer than
at most $nM^2$ possible values. Encode a vector $v$ by
\[
E(v)=\sum_{a=0}^{n-1}v_aB^a,
\qquad T=\{c g^{E(v)}:v\in S\}.
\]
Since $E(v)<B^n\le\operatorname{ord}(\varphi(g))$,
the map from $S$ to $T$ is injective after applying
$\varphi$. Thus $|T|=|S|$.

Let $x,y,z\in S$ encode $\sigma,\tau,\rho$.
Write $a=j-i>0$. The group relation implies
\[
\varphi(g)^{jE(y)}
 =\varphi(g)^{aE(x)+iE(z)}.
\]
Each exponent is the base-$B$ encoding of a vector
whose coordinates lie in $[0,B)$: they are $jy_a$ and
$ax_a+iz_a$, respectively, and $j\le L$. Both
exponents are therefore in $[0,B^n)$.
The order hypothesis makes them equal as integers;
uniqueness of base-$B$ expansion gives
$j y=a x+i z$ coordinatewise. Thus $y$ is a convex
combination of $x$ and $z$ with strictly positive
coefficients. Since all three vectors lie on one sphere,
strict convexity of the squared norm forces $x=z=y$.
In particular $\varphi(\tau/\sigma)=1$.
\end{proof}
This specialization of the formal theorem is
\texttt{exists\_three\_term\_free}
in \leanfile{New/Data/DataExist.lean}{DataExist.lean}.

\section{The class-group and prime choices}
\label{sec:class-group-data}
The class-group step has a short independent proof.
\begin{proposition}
\label{prop:class-pigeonhole}
Let $K$ be a number field with class number $h$.
For $0\le i\le h$, let $I_i,J_i$ be nonzero prime ideals
such that $I_iJ_i$ is principal. Then some distinct
$i,j$ have $I_iJ_j$ principal.
\end{proposition}
\begin{proof}
Two of the $h+1$ classes $[I_i]$ agree. Since
$[I_j][J_j]=1$, the equality $[I_i]=[I_j]$ gives
$[I_iJ_j]=[I_jJ_j]=1$.
\end{proof}
This is \texttt{exists\_mul\_principal} in
\leanfile{New/Data/DataExist.lean}{DataExist.lean}.
In \texttt{exists\_data}, primes $g_0,g_1$ are first
chosen in separated short windows, each in a congruence
class forcing the relevant ideal split. The proposition
crosses two choices to obtain an algebraic integer
$\pi_g$ with $\pi_g\overline{\pi_g}=g_0g_1$.
The large Behrend set is then transferred to the group
of infinite places, giving the thick set.

The record also demands primes in a progression
$1\pmod{M_\omega^2}$ in every logarithmic window of
width $\sqrt{\log b}$ in a specified range. This
supplies the level prime $N$ once its target logarithm
is known. The source packages the analytic prime-in-
progression estimate as a prerequisite to
\texttt{exists\_data}; the present proof of
Proposition~\ref{prop:class-pigeonhole} does not establish
that estimate or the full data-existence theorem.

\section{Rows, columns, and theta lifts}
\label{sec:spaces}
Columns are algebraic forms on a definite unitary group of rank
$k$. Rows are cusp forms on an indefinite unitary group of rank
$r$. The finite group $G$ is a product of general linear groups
over residue fields at chosen places. Its idempotent $e_0$ cuts
out the row space. A Poincar\'e map from columns to rows,
its translates by $G$, and a chain $\Theta_\vartheta$ of $m$
theta lifts provide the forms used in Theorems A--C. The guide
records these definitions, and
\leanfile{New/Trace/GoodIdempotent.lean}{GoodIdempotent.lean}
defines the good row space.

This setup is a large part of the proof. In particular, the
identification of the rows with the representation to which
the height and integral-lattice estimates apply is carried by
the formal development. This manuscript uses the notation
$\Theta_\vartheta$ only for that specific chain.

\section{The row parameter and its weight gap}
\label{sec:row-parameter}
The first indefinite space $Y_0$ has rank $r$.
The row parameter at an infinite place $v$ has entries
\begin{equation}
\label{eq:row-parameter}
\varsigma_v(i)=
\begin{cases}
 k/2-1-i,&i+1<k,\\
 r^8-w_{i+2-k}(v),&i+1\ge k.
\end{cases}
\end{equation}
The source implements this as \texttt{Data.rowParameter}
in \leanfile{New/Seesaw/IndexTypes.lean}%
{IndexTypes.lean}. The first block is the regular
algebraic part; the later entries are pushed down by
the large stage weights. The signs of the diagonal
Hermitian spaces $W_j$ change when the basis number
crosses $k+j-1$. This is \texttt{Data.signsW}.

The gap is quantitative. The declaration
\texttt{Data.varsigma\_add\_w\_gap} in
\leanfile{New/Trace/GoodIdempotent.lean}%
{GoodIdempotent.lean} says that, according to which
side of the sign change an index lies on,
$\varsigma_v(i)+w_j(v)$ is at least $r^8$ or at most
$-r^8$. The gap prevents weights belonging to
different stages from being confused in the theta
and trace calculations. The same file proves
\texttt{isCuspKernel\_paramW}, which puts the
row parameter in the cuspidal range required by the
Poincar\'e and period constructions. The precise
representation-theoretic proof is in the source; the
display \eqref{eq:row-parameter} explains why the
large gap was built into admissibility.

\section{The good projector and the finite group}
\label{sec:good-projector}
The finite group is not an abstract collection of
translations. The source defines
\begin{equation}
\label{eq:finite-group}
G=\prod_{w\in\mathcal W}
\operatorname{GL}_{r}(\mathbb F_w),
\end{equation}
where $\mathcal W$ is a chosen family of places above
the two primes $g_0,g_1$ and $\mathbb F_w$ is the
residue field. The principal-level forms carry a
representation of $G$ by reduction at those places.
For a prescribed family of ``good'' characteristic
polynomial tuples $c$, let $e_c$ be the associated
tuple idempotent. Then
\begin{equation}
\label{eq:good-idempotent}
e_0=\sum_{c\in\mathcal C_{\rm good}}e_c\in\C[G].
\end{equation}
These are \texttt{Data.G},
\texttt{Data.goodTupleFamilies}, and the source's $e_0$
in \leanfile{New/Seesaw/IndexTypes.lean}%
{IndexTypes.lean}. A central-projector lemma in
\leanfile{New/Trace/GoodIdempotent.lean}%
{GoodIdempotent.lean} proves that $e_0$ is a central
self-adjoint idempotent. The row space is
the image of $e_0$ acting on principal-level forms on
$Y_0$; see \texttt{Data.goodForms} and
\texttt{Data.rows} in
\leanfile{New/Trace/GoodIdempotent.lean}%
{GoodIdempotent.lean}.

The projection does two jobs. In Theorem B, a translate
of a column reaches the selected row representation
only after applying $e_0$. In Theorem C, the same
projector lets the good lift decompose into Hecke
eigenspaces whose volumes can be compared. The
reflection-sum lemma for $e_0$ annihilates
specified unwanted reflection sums, a local reason the
projector is useful. None of these facts follows from
idempotence alone: the good-tuple conditions and their
Hecke compatibility are part of the formal source.

\section{Twists and the chain}
\label{sec:twist-data}
A \texttt{Data.TameChar} is a weight-zero character
of $U(1)$ with tame part supported on a finite set
$Q$ of places. The set $Q$ is disjoint from its
complex conjugate and has norm product coprime to
$gN$. A \texttt{Data.Twist} additionally requires
the tame part to be nontrivial at every place of $Q$.
These are exact fields of the structures in
\leanfile{New/Seesaw/IndexTypes.lean}%
{IndexTypes.lean}. The chain requires an
\texttt{MFoldTwist}, so the first $m$ powers of a twist
have the permitted conductor behavior.

The map $\Theta_\vartheta$ is the composition of the
$m$ theta lifts at those powers, restricted to the
rows. Its definition is the source's $\Theta$ in
\leanfile{New/Rows/RowVolume.lean}{RowVolume.lean}.
The trivial twist $1$ supplies the archimedean norm
and covolume bound in Theorem A. The fit twists
$\vartheta_\iota$ are used only in the lower bound:
their altered congruences make the ratio of twisted
and untwisted volumes have a controlled height.

\chapter{Theorem A: short forms}
\label{ch:short}
\section{The two conclusions}
For each $s\in G$, let $B_s$ be a rationally scaled projected
Poincar\'e map from ramified columns to rows. The formal
Theorem A, \texttt{Data.short\_forms}, supplies one positive
rational scaling $q$ and a spanning set $S$ of integral
ramified columns. Its first conclusion bounds every untwisted
lift of every $B_s\phi$, $\phi\in S$, at every complex embedding:
\begin{equation}
\label{eq:short-bound}
\|\Theta_1(B_s\phi)\|_\sigma^2
\le\exp(h_0-mR+L).
\end{equation}
Its second conclusion controls the index of every finite
family selected from the images $B_s(S)$ with $s$ in a finite
set $t\subset G$:
\begin{equation}
\label{eq:a-index}
\delta\le |t|\dim(\mathrm{Column})\,\mathcal B.
\end{equation}
The index is an integral-lattice condition at all relevant
$p$-adic places. It is stronger than a bound on a determinant
over $\C$.

\begin{theorem}[Theorem A in the compact Lean source]
\label{thm:a}
Let $D$ be admissible data. There are $q\in\Q_{>0}$ and
maps $B_s=qP_s$ from ramified columns to rows, where $P_s$
is the projected Poincar\'e map translated by $s\in G$.
There is a subset $S$ of integral ramified columns spanning
the ramified-column space. For every $\phi\in S$, $s\in G$,
and complex embedding $\sigma$, the norm bound
\eqref{eq:short-bound} holds. For every finite $t\subset G$
and every finite family of rows selected from
$\bigcup_{s\in t}B_s(S)$, its index is at most the right
side of \eqref{eq:a-index}.
\end{theorem}
The formal proof is \texttt{Data.short\_forms} in
\leanfile{New/Endgame/ShortAndEnough.lean}{ShortAndEnough.lean}.
Theorem~\ref{thm:a} is stated here at the level of its
mathematical output; the source specifies the index through
\texttt{Data.IndexLe}, including a matrix over each valuation
ring and a bound on the norm of its determinant.

\section{The index is an integral lifting condition}
\label{sec:index-definition}
Here is the exact shape of that condition. For an
independent family $y=(y_1,\ldots,y_n)$, the assertion
\texttt{IndexLe}$(y,\delta)$ means that there are real
numbers $e_\ell$ indexed by rational primes, with
\begin{equation}
\label{eq:index-budget}
\sum_{\ell\in U}e_\ell\le\delta
\quad\text{for every finite set of primes }U,
\end{equation}
such that the following holds. At every $\ell$-adic
embedding $\iota$ and every twist $\vartheta$ of conductor
prime to $\ell$, there is a matrix $u$ over the valuation
ring $\mathcal O_\iota$ with
\begin{equation}
\label{eq:index-matrix}
|\det u|_\iota\ge e^{-e_\ell},\qquad
\Theta_\vartheta\!\left(\sum_j u_{ij}y_j\right)
\text{ integral at }\iota\quad\text{for every }i.
\end{equation}
This is \texttt{Data.IndexLe} in
\leanfile{New/Endgame/ShortAndEnough.lean}%
{ShortAndEnough.lean}. The definition quantifies over
all finite sets in \eqref{eq:index-budget}; it does not
assume in advance that $\sum_\ell e_\ell$ converges.
The matrix may depend on the embedding and the twist.
Its determinant, not an entrywise bound, is what later
costs $2\delta$ in the Gram determinant.

The rational scaling in Theorem~\ref{thm:a} is elementary.
\begin{proposition}
\label{prop:rational-scaling}
Let $T$ be a finite set of rational primes and let
$g_\ell\in\Z$ vanish outside $T$. There is a positive
$q\in\Q$ with
\[
\log q=\sum_{\ell\in T}g_\ell\log\ell,
\qquad |q|_p=|p|_p^{g_p}
\quad\text{for every rational prime }p.
\]
\end{proposition}
\begin{proof}
Take $q=\prod_{\ell\in T}\ell^{g_\ell}$.
The real logarithm is additive, and unique factorization
gives $v_p(q)=g_p$. With $|p|_p=p^{-1}$, the second
identity follows.
\end{proof}
This is \texttt{exists\_rat\_log\_eq\_and\_padic\_norm\_eq}
in \leanfile{New/Endgame/ShortAndEnough.lean}%
{ShortAndEnough.lean}. The source chooses
$g_\ell=\lfloor\mathrm{allowedLoss}(\ell)-
\mathbf1_{\ell\in\mathcal P_{\rm ns}}
\mathrm{gain}(\ell)\rfloor$ on the nonsplit primes
and the level prime. Its real logarithm offsets the
positive gain term in \eqref{eq:weyl-bound}; its local
norms shift the integrality condition by the same integer
with the opposite valuation sign.

% Superseded by Proposition~\ref{prop:rational-scaling} above.
% The scale q is a product of rational prime powers; its real and
% p-adic norms encode gains and losses with opposite signs.
The deeper work in Theorem A lies in the projected Poincar\'e
series, the seesaw identity, the Weyl norm estimate, and the
local divisibility of theta lifts. The exact import chain is
visible at the head of
\leanfile{New/Endgame/ShortAndEnough.lean}{ShortAndEnough.lean}.

\section{The two estimates behind Theorem A}
\label{sec:a-mechanisms}
The seesaw identity compares two constructions of the same
chain of forms. There is a nonzero scalar $c$ such that for
each tame twist $\vartheta$ a good finite datum $\Psi$ has
\begin{equation}
\label{eq:seesaw}
\iota_\vartheta\bigl(\mathcal P(c\phi)\bigr)
 =B^\Psi_\vartheta\phi
\qquad(\phi\text{ a ramified column}).
\end{equation}
Here $\mathcal P$ is a Poincar\'e series, $\iota_\vartheta$
the $m$-stage theta chain through the indefinite spaces, and
$B^\Psi_\vartheta$ its definite-space expansion. This is
\texttt{seesaw} in
\leanfile{New/Seesaw/SeesawTheorem.lean}{SeesawTheorem.lean}.
Equation~\eqref{eq:seesaw} lets the construction control
the norm of the row lift through columns, whose local
divisibility can be estimated.

\section{The finite-cut seesaw induction}
\label{sec:finite-cut}
The seesaw identity is not a formal consequence of the
usual one-step theta correspondence. At stage $j$ the
source unfolds a period on an indefinite space into a
period on a definite space. The unfolding initially
involves a quotient by a compact open subgroup; the proof
must replace it by a finite sum whose denominators remain
controlled at every prime away from the twist conductor.

For a stage $j$ and a twist $\vartheta$, let
$\mathcal J_j(\vartheta)$ denote the level subgroup in
\texttt{Data.levelY}. The declaration
\texttt{Data.exists\_stage\_cut} in
\leanfile{New/Seesaw/SeesawTheorem.lean}%
{SeesawTheorem.lean} supplies a compact open $A_j$
containing the column level and a finite set
$T_j(\vartheta)$ of cosets modulo
$\mathcal J_j(\vartheta)$. It is stable under the
relevant Bessel subgroup inside $A_j$. Representatives
are integral away from the conductor of $\vartheta$;
at stage zero they lie in the chosen level. The
covering multiplicity of each coset that contributes
to the unfolded period divides a power of the conductor.
Consequently the reciprocals of those multiplicities
have no denominator away from that conductor.

The next declaration,
\texttt{Data.exists\_isAdmissibleComb\_unfolding},
builds the finite sums by downward induction on $j$.
Its induction variable is an admissible linear combination
of values of a form on the finite coset set. At the
last stage, the integral is evaluated by a finite
unfolding. Moving from $j+1$ to $j$ replaces one period
by a sum over $T_j(\vartheta)$ weighted by inverse
covering multiplicities. The support and integrality
properties from \texttt{exists\_stage\_cut} show that
the new combination remains admissible. At $j=0$ the
period is the Poincar\'e series pairing. The same
induction tracks the level action and a scalar independent
of the input column.

Finally \texttt{Data.seesaw} chooses a nonzero $c$ and
a good datum $\Psi$ from that invariant combination.
This gives \eqref{eq:seesaw} for every ramified column
and every tame twist. The paragraph above is the
construction route in the pinned source. The existence
of the cuts, the period identities, and the invariant
combination are used as formal results here; a full
conventional proof would need their local measure and
denominator normalizations.

The archimedean estimate
\path{weylNormSq_Bmu_le} in
\leanfile{New/Columns/LocalStatement.lean}{LocalStatement.lean}
has the following form. If the conjugate norm square of a
column is at most $e^B$ and $\Psi$ is good, then the norm
square of its $B^\Psi$ image is at most
\begin{equation}
\label{eq:weyl-bound}
\exp\!\left(B+\mathrm{heightPerRow}
  +2\sum_{p\in\mathcal P_{\rm ns}}
       \mathrm{gain}(p)\log p+\mathcal B\right).
\end{equation}
The finite set $\mathcal P_{\rm ns}$ consists of selected
small nonsplit primes. A rational factor $q$ is then chosen
so its logarithm cancels the gains in \eqref{eq:weyl-bound},
up to the allowed local losses. The bound
\eqref{eq:height-per-row} turns the result into
\eqref{eq:short-bound}.

At a prime $p$, the local statement
\texttt{exists\_indexDvd\_Bmu\_divisible} supplies a
submodule of columns of controlled index. For each integer
\[
\beta<\mathbf1_{p\in\mathcal P_{\rm ns}}
       \mathrm{gain}(p)-\mathrm{allowedLoss}(p)+1,
\]
the appropriately $p^{-\beta}$-scaled $B^\Psi$ image of
its ramified elements is integral at the relevant $p$-adic
place. The index exponent is at most
$\mathrm{allowedLoss}(p)\dim(\mathrm{Column})$.
This is the local origin of \eqref{eq:a-index}; summing
the local losses and choosing $q$ give the global bound.
The source treats collisions at small nonsplit primes and
the level prime separately. Both cases are part of the
stated theorem, and neither follows merely from
\eqref{eq:seesaw}.

\section{Marked Hecke operators and local divisibility}
\label{sec:marked-hecke}
The local work is concentrated at a small nonsplit prime
$p<x_1$ which can pair two top places. The set $\Lambda$
of marked places above $p$ records where a Hecke operator
is applied. If $Y_\Lambda$ is the product of those
operators, the source's
\texttt{Data.IsMarked.hasType\_prod\_markOperator\_pow}
shows that $Y_\Lambda^m$ preserves the integral-column
lattice and that its image, after adjoining a prescribed
$p$-power multiple of all columns, has index dividing
\begin{equation}
\label{eq:marked-index}
p^{3\,\mathrm{loss}(p)\dim(\mathrm{Column})}.
\end{equation}
The exponent is a cost of making the columns satisfy the
divisibility condition; it does not enter the archimedean
norm of each selected column directly.

Under the stronger \texttt{IsFullMark} condition, the
source proves that for an integral column $v$ the image
$B^\Psi_\vartheta(Y_\Lambda^m v)$ has $p$-adic
valuation at least
\begin{equation}
\label{eq:mark-gain}
\mathrm{gain}(p)-2\,\mathrm{loss}(p).
\end{equation}
The local inequality is
\path{Data.IsGoodDatum.gaugeLe_Bmu_of_isMarked}.
It is in \leanfile{New/Columns/LocalStatement.lean}%
{LocalStatement.lean}. A second route treats collisions:
without a full mark, \texttt{gaugeLe\_collisionSum}
allows an extra $m\,\mathrm{collisionSum}(p,\varnothing)$
loss, and \texttt{gaugeLe\_pow\_jetCeil\_smul}
repairs it by multiplying the input by
$p^{\mathrm{jetCeil}(p)}$. The local index theorem combines
these routes and handles the exceptional level prime
separately. Its final output is the submodule and integer
$\beta$ described after \eqref{eq:a-index}, namely
\path{Data.exists_indexDvd_Bmu_divisible}.

The proof of \eqref{eq:mark-gain} has a useful accounting
identity. At each of the $m$ stages, moving the marked
operator through the theta kernel contributes a power
of $p$ from a local Hecke relation. A \texttt{liftLoss}
term measures the possible loss from the low-column
and jet exponents. Summing over stages and marked places
gives the initial gauge exponent in
\texttt{gaugeLe\_markOperator\_pow}. The full-mark
condition bounds that sum by the two losses in
\eqref{eq:mark-gain}. The Hecke transfer itself is
\texttt{exists\_hecke\_transfer}; it writes a large Hecke
operator after a lift as an integral combination of
smaller ones. Its representation and local-kernel
identities are formal inputs to this account.

\section{The archimedean estimate as a chain bound}
\label{sec:arch-chain}
The source proves \eqref{eq:weyl-bound} by expanding
$B^\Psi_1\phi$ into finitely many chains of theta
kernels. The datum $\Psi$ has controlled support and
denominators because it came from the finite-cut
construction. A jet-norm estimate bounds each stage;
the sum of the main jet terms is
\texttt{Data.heightPerRow}. The small nonsplit primes
contribute twice their gain to the logarithm of the
squared norm. The remaining combinatorial and rounding
loss is $\mathcal B$. This is the route through
\texttt{exists\_chain\_eq\_sum} and
\texttt{weylNormSq\_Bmu\_le} in
\leanfile{New/Columns/LocalStatement.lean}%
{LocalStatement.lean}.

For the columns selected in Theorem A, the initial
conjugate squared norm is bounded by
$\exp(2\log q+\mathcal B)$. Inserting this in
\eqref{eq:weyl-bound} and using
\eqref{eq:height-per-row} gives
\[
\log\|\Theta_1(B_s\phi)\|_\sigma^2
\le h_0-mR+3\mathcal B
 +2\sum_{p\in\mathcal P_{\rm ns}}
      \mathrm{gain}(p)\log p+2\log q.
\]
The exponent of each prime in $q$ is the floor of its
allowed loss minus its gain. Hence the last two terms
are bounded by the sum of allowed local losses, plus
rounding. The source's
\texttt{sum\_allowedLoss\_mul\_log\_le} and scale bounds
absorb this in $L-3\mathcal B$, yielding
\eqref{eq:short-bound}. This is an accounting derivation
conditional on the archimedean and local source lemmas;
it does not prove their analytic estimates.

\section{Why the level matters}
Our notation $h_0=m k\mathcal U$ denotes the source's
\texttt{Data.mainHeight}; $\mathcal U$ is the cost of one
direction. The source proves
\begin{equation}
\label{eq:height-per-row}
\mathrm{heightPerRow}\le h_0-mR+\mathcal B.
\end{equation}
This is \texttt{Data.heightPerRow\_le}. Its proof compares a
factorial bound for the size of jets with the chosen logarithm
of the level prime $N$. The leading terms cancel. The
remainder $-mR$ is therefore visible in the short-form
estimate. Choosing $N$ after the relevant scales would not
give this cancellation automatically; the short level window
is part of admissibility.

Write $\bar w=f^{-1}\sum_\sigma\log w_{\rm nom}(\sigma)$
and $G=\log(g\,\mathrm{rd}_E/(2\pi))$, where
$\mathrm{rd}_E$ is the root discriminant. The source's
main ask is
\begin{equation}
\label{eq:main-ask}
A_{\rm main}=a(\log a-1-\operatorname{CS}(d)-\bar w).
\end{equation}
Its level target is
\begin{equation}
\label{eq:level-target}
\log N\ge m(\bar w+G)-G-\tfrac{15}{2}\log b.
\end{equation}
The difference between \texttt{jetSize}$(a)$ and
$A_{\rm main}$ satisfies
\begin{equation}
\label{eq:jet-excess}
\mathrm{jetSize}(a)-A_{\rm main}
\le a(\bar w+G)+1+\tfrac12\log a.
\end{equation}
The proof of \eqref{eq:jet-excess} substitutes the
definition of the period $\Omega$ and applies the
upper Stirling bound
$\log(a!)\le a\log a-a+1+\tfrac12\log a$.

\begin{proposition}
\label{prop:level-cancellation}
Let $D$ be admissible data. Assume the level bound
\eqref{eq:level-target}, the jet bound
\eqref{eq:jet-excess}, and the source's scale estimate
\[
\#\mathrm{tops}\left(m(1+\tfrac12\log a)
 +a(G+\tfrac{15}{2}\log b)\right)\le\mathcal B.
\]
Then $\mathrm{heightPerRow}\le h_0-mR+\mathcal B$.
\end{proposition}
\begin{proof}
Unfolding the height and the definition of the remainder
in \texttt{ShortAndEnough.lean} gives the exact identity
\[
\mathrm{heightPerRow}=h_0-mR+
\#\mathrm{tops}\bigl(m(\mathrm{jetSize}(a)-A_{\rm main})
 -a\log N\bigr).
\]
Insert \eqref{eq:jet-excess} and
\eqref{eq:level-target}. The $ma(\bar w+G)$ terms
cancel. The displayed scale estimate bounds what remains
by $\mathcal B$.
\end{proof}
This is the mathematical calculation in
\texttt{Data.heightPerRow\_le} of
\leanfile{New/Endgame/ShortAndEnough.lean}%
{ShortAndEnough.lean}. The existence of $N$ in the
specified short interval and the scale estimate are
separate parts of the admissible-data construction.

\chapter{Theorem B: enough forms}
\label{ch:enough}
\section{Almost spanning with few translates}
The formal Theorem B, \texttt{Data.enough\_forms}, gives a
finite set $t\subset G$ such that
\begin{equation}
\label{eq:few-translates}
|t|\dim(\mathrm{Column})
\le5D_{\rm row}\log(1+b^2)
\end{equation}
and the sum of the images of the maps $B_s$, $s\in t$,
has dimension at least $(1-4/b^2)D_{\rm row}$.
This is a dimension statement, not the assertion that every
row is a projected Poincar\'e series.

\begin{theorem}[Theorem B in the compact Lean source]
\label{thm:b}
Let $D$ be admissible data with column space $V$, row
space $W$, and projected Poincar\'e maps $P_s:V\to W$
for $s\in G$. Both $V$ and $W$ are nonzero. There is a
finite $t\subset G$ satisfying \eqref{eq:few-translates}
and
\[
\dim\sum_{s\in t}\operatorname{im}P_s
 \ge(1-4/b^2)\dim W.
\]
\end{theorem}
The formal proof is \texttt{Data.enough\_forms} in
\leanfile{New/Endgame/ShortAndEnough.lean}{ShortAndEnough.lean}.

Its proof invokes \path{exists_few_translates} after
\path{ramifiedColumns_fill_rows}. The latter is the
representation-theoretic filling statement. In the combination
with Theorem A, choose a basis from the translated images.
Equation~\eqref{eq:few-translates} inserted into
\eqref{eq:a-index} gives \eqref{eq:index}; the uniform
norm bound \eqref{eq:short-bound} yields the covolume bound
\eqref{eq:upper-covolume}. This is the construction of the
family used in Chapter~\ref{ch:assembly}.

\section{Why the Poincar\'e translates spread}
\label{sec:filling}
The map before the final rational scaling has the form
\[
P_s=\operatorname{pr}_{e_0}\,\rho(s)
   \mathcal P\,\iota_{\rm ram},\qquad s\in G.
\]
Here $\mathcal P$ is the Poincar\'e series, $\rho$ is
the finite-level group action, $e_0$ is the central
idempotent selecting good rows, and $\iota_{\rm ram}$
is the ramifying operator on columns. The source
normalizes the projection by $|G|$. Its theorem
\texttt{Data.ramifiedColumns\_fill\_rows} states the
spread condition \eqref{eq:spread} for these maps with
$\eta=b^{-2}$; see
\leanfile{New/Filling/FillingTheorem.lean}%
{FillingTheorem.lean}. This is the deep part of Theorem B.

The proof begins with a Petersson pairing for the
Poincar\'e series. Matrix coefficients of translates
of columns become the \texttt{amplitudeSum} of a
finite-level model. The source estimates its second
moment. The main energy is concentrated on a Bessel
subgroup; \texttt{energy\_compl\_besselS\_le}
bounds the complementary energy by $b^{-12}$ times
the Bessel energy. The fixing-set estimate
\texttt{energy\_fixingSet\_le} controls the mass of
terms with unwanted stabilizers. These estimates are
proved in \leanfile{New/Filling/Moments.lean}%
{Moments.lean}; they depend on the admissible data's
three-term-free thick set and the rank gaps.

The linear-algebra conversion is worth making explicit.
For a map $P:\C^n\to H$, a finite family
$v_i=Pe_i$, and an orthonormal basis $b$ of $H$, set
\[
\mathcal D(v)=\sum_i\overline{v_i}^{\,b}\otimes v_i.
\]
The source proves
\begin{equation}
\label{eq:dyad-pairing}
\langle\mathcal D(v),\mathcal D(w)\rangle
 =\sum_{i,j}|\langle w_j,v_i\rangle|^2.
\end{equation}
For a subspace $A\subset H$, choosing $w$ as an
orthonormal basis of $A$ makes the right side the
Hilbert--Schmidt square of the projection of $P$ to
$A$. This is \texttt{inner\_dyadSum} in
\leanfile{New/Heights/OneDeterminant.lean}%
{OneDeterminant.lean}. A trace comparison in
\texttt{finrank\_sub\_le\_re\_inner\_dyadSum}
then turns closeness of $P^*P$ to the identity and a
small bad component into a lower bound for the image
outside any prescribed subspace. Averaging over $s\in G$
shows that at least one translate has the required
rank outside $A$.

For finite-dimensional $V,W$ with $W\ne0$ and maps
$B_s:V\to W$, the source calls the family
$\eta$-\emph{spread} if, for every subspace $F\subseteq W$,
some $s$ satisfies
\begin{equation}
\label{eq:spread}
\dim B_s^{-1}F\le
\left(\frac{1+\dim F/\dim W}{2}+\eta\right)\dim V.
\end{equation}
This is \texttt{IsSpread} in
\leanfile{New/Heights/OneDeterminant.lean}%
{OneDeterminant.lean}.

The abstract output of the second-moment argument is
\path{Fits.exists_translate_finrank_comap_le}
in \texttt{Moments.lean}. Before ramification it yields
an \texttt{IsSpread} error $(19/2)b^{-4}$. The final
ramifying operator and the controlled bad share consume
part of this margin. The scale inequalities enlarge the
error to $b^{-2}$, giving
\texttt{ramifiedColumns\_fill\_rows}. These constants
are not obtained from the elementary greedy lemma below;
their proof depends on the automorphic Petersson and
finite-level energy identities. This manuscript takes
those identities at their formal source locators.

\begin{proposition}
\label{prop:rows-nonzero}
Let $V$ be a nonzero finite-dimensional complex vector
space, let $W$ be finite-dimensional, and suppose a
family $B_s:V\to W$ is spread with $0<\eta<1/2$.
Then $W$ is nonzero.
\end{proposition}
\begin{proof}
If $W=0$, the only subspace is $F=0$. The spread
inequality says $\dim V\le(1/2+\eta)\dim V$, contrary
to $\dim V>0$.
\end{proof}
The source uses this argument as
\texttt{Data.finrank\_rows\_ne\_zero}.

\section{A finite-dimensional spreading lemma}
\label{sec:spread}
The quantitative linear algebra is separated from the
automorphic filling assertion. Inequality~\eqref{eq:spread}
measures how much of a new image can lie in a subspace
already reached.

\begin{theorem}
\label{thm:few-translates}
Let $V$ and $W$ be finite-dimensional complex vector spaces,
with $V$ and $W$ nonzero. Let $(B_s:V\to W)_{s\in\Theta}$ satisfy
\eqref{eq:spread} for some $\eta>0$. There is a finite set
$t\subset\Theta$ of indices such that
\[
|t|\dim V\le5\dim W\log(1+1/\eta)
\]
and
\[
\dim\sum_{s\in t}\operatorname{im}B_s
 \ge(1-4\eta)\dim W.
\]
\end{theorem}
\begin{proof}
Put $d=\dim V$ and $D=\dim W$. If $\eta\ge1/4$, take
$t=\varnothing$. Suppose $0<\eta<1/4$. Apply the spread
condition with $F=0$. Some $B_s$ has rank at least
$(1/2-\eta)d>d/4$, so $d<4D$.

Start with $F_0=0$ and write $q_i=\dim(W/F_i)$. Whenever
$q_i>4\eta D$, choose $s$ from \eqref{eq:spread} and set
$F_{i+1}=F_i+\operatorname{im}B_s$. Rank-nullity gives
\[
\dim(F_{i+1}/F_i)
=d-\dim B_s^{-1}F_i
\ge\left(\frac{q_i}{2D}-\eta\right)d
\ge\frac{q_i d}{4D}.
\]
Thus $q_{i+1}\le q_i(1-d/(4D))$. The chosen $s$ is new:
an old image lies in $F_i$ and would add zero dimension.
After at most
\[
k=\left\lceil\frac{4D}{d}\log\frac1{4\eta}\right\rceil
\]
steps, either the construction has stopped or
$q_k\le D e^{-kd/(4D)}\le4\eta D$. Also
\[
kd\le4D\log\frac1{4\eta}+d
  <4D\left(\log\frac1{4\eta}+1\right)
  \le5D\log(1+1/\eta).
\]
For the last inequality, use $\log4>1$ and
$\log(1+1/\eta)\ge\log(1/\eta)$. The selected
indices form the required set $t$.
\end{proof}
The source proves this as \texttt{exists\_few\_translates}.
An alternative proof in the source uses the potential
\[
\Phi(F)=5\dim W\log\left(1+
\frac{\dim(W/F)}{\eta\dim W}\right).
\]
Starting with $F=0$, add an image $B_s(V)$ when the spread
condition makes the quotient dimension drop. The proof shows
that the decrease in $\Phi$ pays at least $\dim V$ for each
added translate. A maximal set under the resulting budget
has quotient dimension at most $4\eta\dim W$, which gives
the theorem. In Theorem B the source takes $\eta=b^{-2}$.
The remaining burden is to prove \eqref{eq:spread} for the
specific Poincar\'e maps; that is the role of the filling
modules and is not a consequence of this linear-algebra lemma.

\chapter{Theorem C: integral forms are long}
\label{ch:long}
\section{Statement of the lower bound}
\begin{theorem}[Theorem C in the compact Lean source]
\label{thm:c}
Let $D$ be admissible data with row dimension
$D_{\rm row}>0$. Let $y_1,\ldots,y_n$ be independent rows,
where $n\ge(1-4/b^2)D_{\rm row}$. Assume their index is
at most $\delta$ in the sense of \texttt{Data.IndexLe},
and for every complex embedding $\sigma$ their untwisted
theta lifts have squared covolume at most $e^V$. Then
\[
D_{\rm row}(h_0-3L)\le V+2\delta.
\]
\end{theorem}
The formal proof is \texttt{Data.integral\_forms\_long} in
\leanfile{New/Endgame/LongForms.lean}{LongForms.lean}.

\section{The height of a family indexed by primes}
\label{sec:height-definition}
The height in Theorem C is not the usual logarithmic
height of a single algebraic number. A family has one
complex value $z_\iota$ for every rational prime $\ell$
and every field isomorphism
$\iota:\C\simeq\overline{\Q}_\ell$.
The definition \texttt{HeightLe} in
\leanfile{New/Heights/Heights.lean}{Heights.lean}
says that $\operatorname{HeightLe}(c,z)$ holds when there
are a nonzero $H\in\C$, a real $\alpha$, and real
numbers $e_\ell$ such that
\begin{align}
\label{eq:height-arch}
|\sigma H|&\le e^\alpha
       &&(\sigma\in\operatorname{Aut}\C),\\
\label{eq:height-finite}
|\iota(H z_\iota)|_\ell&\le e^{e_\ell}
       &&(\iota:\C\simeq\overline\Q_\ell),
\end{align}
and, for every sufficiently large cutoff $M$,
\begin{equation}
\label{eq:height-budget}
\alpha+\sum_{\substack{\ell<M\\\ell\text{ prime}}}
 e_\ell\le c.
\end{equation}
The common factor $H$ clears denominators. It can be
chosen differently for different families, but not
separately for each $\ell$. This distinction is needed
when multiplying height bounds.

\begin{proposition}
\label{prop:height-product}
Let $z,z'$ be families indexed by prime-adic
isomorphisms. If $\operatorname{HeightLe}(c,z)$ and
$\operatorname{HeightLe}(c',z')$, then
$\operatorname{HeightLe}(c+c',zz')$.
\end{proposition}
\begin{proof}
Multiply the two common factors $H,H'$. Add their
archimedean exponents and add the local exponents
$e_\ell,e'_\ell$. The inequalities
\eqref{eq:height-arch}--\eqref{eq:height-budget}
multiply and add respectively, giving the asserted
budget for every sufficiently large cutoff.
\end{proof}
This is \texttt{HeightLe.mul}. Iteration gives the
product and power rules used in the volume factorization.

The terminal principle is \texttt{HeightLe.nonneg}:
\begin{equation}
\label{eq:height-one}
\operatorname{HeightLe}(c,1)\quad\Longrightarrow\quad c\ge0.
\end{equation}
For a witness $H$, the source first obtains the
algebraicity needed to pass to the number field $\Q(H)$.
It then applies the product formula to its norm, comparing
the average archimedean logarithm with the $\ell$-adic
logarithms. The finite-cutoff form
\eqref{eq:height-budget} avoids assuming absolute
convergence of $\sum e_\ell$. The algebraicity step and
normalization of local norms are part of the formal proof
in \texttt{Heights.lean}; this manuscript uses
\eqref{eq:height-one} as an imported result.

\section{Volume ratios and heights}
The compact proof attaches a volume $\mathfrak z(\vartheta)$
to a twist of the row representation. It factors into a scalar
and the volumes of the individual theta lifts. The factorization
is \texttt{volume\_iota\_eq\_prod\_liftVolume} in
\leanfile{New/Rows/RowVolume.lean}{RowVolume.lean}.
For each $p$-adic embedding $\iota$, the proof selects a fit
twist $\vartheta_\iota$. Its congruence properties are
\texttt{IsFitFamily}; existence is
\texttt{exists\_isFitFamily} in
\leanfile{New/Selmer/FitTwistsExist.lean}{FitTwistsExist.lean}.

The fit family is chosen \emph{separately} at each
$\ell$-adic embedding. There need not be a single global
twist good at every prime. The declaration
\texttt{exists\_mFoldTwist\_selmerLengthLe} in
\leanfile{New/Selmer/TwistForRows.lean}%
{TwistForRows.lean} produces a twist whose first $m$
powers have controlled Selmer lengths for the finitely
many row Galois representations under consideration.
It also bounds the valuations of ratios attached to
ramified pairs. The twisting characters are chosen
anticyclotomically, with tame generators at auxiliary
places and a bounded sum of local deviations from $1$.
The source's \texttt{IsAuxTwist} records those conditions.

The denominator theorem
\path{denominatorLe_rowBudget} turns a short Selmer
length into
an $\ell$-adic bound for projection onto the lift image.
Its source is
\leanfile{New/Selmer/DenominatorBound.lean}%
{DenominatorBound.lean}.
The statement requires a cut spanned by forms whose
Galois representations satisfy the row condition and
whose twisted Selmer groups have the specified finite
length. It gives a denominator cost
$r\,\mathrm{rowBudget}(D,\ell)\log\ell$.
The construction of \texttt{exists\_isFitFamily}
then assembles these choices and proves
\begin{equation}
\label{eq:fit-budget}
\sum_{\ell\in U}
 \mathrm{rowBudget}(D,\ell)\log\ell\le B_0
\quad\text{for every finite set of primes }U,
\end{equation}
with a separate $x^{999}$ budget for crude scalar
factors. The formal proof uses Galois representations,
Selmer groups, patching, and a factorization bound.
Those modules are the deepest arithmetic inputs to
Theorem C. This manuscript records their exact role
and does not derive the Selmer bound conventionally.

The proof uses two height bounds. The first controls the
squared ratio $\mathfrak z(1)^2/
\mathfrak z(\vartheta_\iota)^2$ by
\begin{equation}
\label{eq:ratio-one}
2D_{\rm row}(2L-h_0).
\end{equation}
The second bounds the height of
$\mathfrak z(\vartheta_\iota)/\mathfrak z(1)$ by
\begin{equation}
\label{eq:ratio-two}
V+2\delta+D_{\rm row}b.
\end{equation}
These are \texttt{heightLe\_volume\_ratio\_sq} and
\texttt{heightLe\_volume\_rows} in
\leanfile{New/Endgame/LongForms.lean}{LongForms.lean}.
The first follows from the lift-volume height gain and bounds
for scalar factors. The second extends an almost spanning
family to a basis with cruder rows and applies a determinant
estimate. Its hypotheses include linear independence, the
near-spanning bound, index control, and the covolume bound at
every complex embedding.

The volume factorization has a precise shape:
\begin{equation}
\label{eq:volume-factorization}
\mathfrak z(\vartheta)
=|c(\vartheta)|^{D_{\rm row}}
 \prod_{j=1}^m\lambda_j(\vartheta^j),
\end{equation}
where $\lambda_j$ is the squared determinant volume of the
$j$th good lift. The scalar $c(\vartheta)$ carries a
controlled square cost. This is
\texttt{volume\_iota\_eq\_prod\_liftVolume} in
\leanfile{New/Rows/RowVolume.lean}{RowVolume.lean}.
It reduces the height of a whole chain to the sum of
the heights of the individual lift-volume ratios and
one scalar correction.

For a fit family $\vartheta_\iota$ and $1\le j\le m$,
\texttt{heightLe\_liftVolume\_ratio} in
\leanfile{New/Rows/LiftVolumeHeight.lean}%
{LiftVolumeHeight.lean} bounds the height of
$\lambda_j(1)/\lambda_j(\vartheta_\iota^j)$ by
\begin{equation}
\label{eq:per-lift-height}
D_{\rm row}k\bigl((\log b)B_0-\mathcal U\bigr),
\qquad B_0=\frac{\mathcal B}{mk}.
\end{equation}
Summing \eqref{eq:per-lift-height} over the $m$ lifts gives
$D_{\rm row}(L-h_0)$, since $h_0=mk\mathcal U$.
The scalar in \eqref{eq:volume-factorization} accounts for
the remaining error in \eqref{eq:ratio-one}. This is how
the row chain repays the cost of the idle directions in the
height calculation.

\section{Comparison forms and compressed Gram matrices}
\label{sec:compressed-gram}
The bound \eqref{eq:per-lift-height} is proved one lift at
a time. Fix $j$ and a twist $\vartheta$. Write
$u_\vartheta$ for the $j$th good theta lift and
$L_\vartheta=\operatorname{im}u_\vartheta$.
The declaration \texttt{exists\_comparison\_forms} in
\leanfile{New/Rows/ComparisonForms.lean}%
{ComparisonForms.lean} supplies a table space, a
surjection $A$ from tables onto good forms, comparison
maps $\Xi_\vartheta$ into the lift's target, and nonzero
scalars $\gamma_\vartheta$ satisfying
\begin{equation}
\label{eq:comparison-adjoint}
u_\vartheta^*\Xi_\vartheta
 =\gamma_\vartheta A.
\end{equation}
The squared ratio
$|\gamma_\vartheta/\gamma_1|^2$ is a ratio of the
source's \texttt{heckeValueProd}. The comparison maps
have an archimedean squared-norm bound at the trivial
twist. At primes not dividing the conductor, they take
integral tables to integral forms. Their images satisfy
the testing condition needed to bound the determinant
on the good-lift subspace.

For a nonzero scalar $h$, a map $\Xi:T\to V$, and a
subspace $L\subset V$, define the compressed Gram
operator on $L$ by
\begin{equation}
\label{eq:compressed-gram}
C_L=h\,\operatorname{pr}_L\Xi\Xi^*|_L.
\end{equation}
This operator gathers the table comparisons into one
determinant. Its elementary identity is the following.
\begin{proposition}
\label{prop:compressed-det}
Let $u:U\to V$ be an injective map of finite-dimensional
complex inner-product spaces, with $L=\operatorname{im}u$.
For $C_L$ as in \eqref{eq:compressed-gram},
\[
\operatorname{vol}(u)^2\det C_L
 =h^{\dim L}\operatorname{vol}(\Xi^*u)^2,
\]
where $\operatorname{vol}(v)^2=\det(v^*v)$.
\end{proposition}
\begin{proof}
Choose orthonormal bases of $U$ and $L$. The matrix
of $u:U\to L$ is invertible, say $U_0$. The matrix
of $C_L$ is $hXX^*$, where $X$ is the matrix of
$\operatorname{pr}_L\Xi$. Hence the left side is
$|\det U_0|^2h^{\dim L}\det(XX^*)$. The Gram matrix
of $\Xi^*u$ is $U_0^*XX^*U_0$, whose determinant is
the same expression without $h^{\dim L}$.
\end{proof}
This is \texttt{det\_compressedGramOp} in
\leanfile{New/Rows/LiftVolumeHeight.lean}%
{LiftVolumeHeight.lean}. Equation
\eqref{eq:comparison-adjoint} replaces
$\Xi_\vartheta^*u_\vartheta$ by
$\overline{\gamma_\vartheta}A^*$, so the dependence
on the twist is concentrated in $\gamma_\vartheta$
and $\det C_{L_\vartheta}$.

The same file bounds the determinant at complex
embeddings using the size of $\Xi_1$; its declaration
is \path{norm_det_compressedGramOp_arch_le}.
At finite places, integrality and the projection
denominator bound give
\path{norm_det_compressedGramOp_padic_le}.
In the local bound, if the denominator
parameter is $\beta$, then
\begin{equation}
\label{eq:compressed-local}
|\det C_{L_\vartheta}|_\ell
 \le\exp(\dim L_\vartheta\, r\beta).
\end{equation}
The archimedean counterpart uses the norm bound for
comparison forms and a positivity argument on the
eigenvalues of $C_{L_1}$. The fit-twist budget
\eqref{eq:fit-budget} sums the local exponents.
The Hecke-value ratio contributes a separate height
bound in \texttt{heightLe\_inv\_heckeValueProd\_div}
of \leanfile{New/Rows/HeckeValueHeight.lean}%
{HeckeValueHeight.lean}. These estimates, inserted in
Proposition~\ref{prop:compressed-det}, give
\eqref{eq:per-lift-height}. The source proves all
integrality, testing, and eigenvalue hypotheses; the
determinant identity alone does not supply them.

\section{The product-formula comparison}
Multiply the square of the ratio in
\eqref{eq:ratio-two} by the ratio in
\eqref{eq:ratio-one}. The resulting family is identically
$1$. Its height is nonnegative. Thus the two upper bounds
cannot have a negative sum. After absorbing
$D_{\rm row}b$ into the allowed error, the source obtains
\eqref{eq:lower-covolume}. This is the last bound needed
for Proposition~\ref{prop:comparison}.

\begin{proposition}
\label{prop:height-comparison}
Let $D$ be admissible data, and suppose the height
bounds \eqref{eq:ratio-one} and \eqref{eq:ratio-two}
hold for a fit family. If $b\le L$, then
$D_{\rm row}(h_0-3L)\le V+2\delta$.
\end{proposition}
\begin{proof}
For each $\iota$, put
$A_\iota=\mathfrak z(\vartheta_\iota)/\mathfrak z(1)$
and $B_\iota=\mathfrak z(1)^2/
\mathfrak z(\vartheta_\iota)^2$. The volumes are nonzero
by \texttt{volume\_iota\_eq\_prod\_liftVolume}, and
$A_\iota^2B_\iota=1$. Proposition~\ref{prop:height-product}
and the power rule give
\[
\operatorname{HeightLe}\!\left(
  2(V+2\delta+D_{\rm row}b)
   +2D_{\rm row}(2L-h_0),1\right).
\]
By \eqref{eq:height-one}, the first argument is
nonnegative. Dividing by $2$ and rearranging gives
$D_{\rm row}(h_0-2L-b)\le V+2\delta$.
Finally $b\le L$ gives the claim.
\end{proof}
This calculation is the last part of
\texttt{Data.integral\_forms\_long}. The hypothesis
$b\le L$ is one of the source's explicit scale
comparisons, discharged by \texttt{D.le\_of\_hasLogBound}.

The product formula is the short final step; the substantive
dependencies are the fit twists, the factorization of volumes,
the local and global height estimates, and the determinant
estimate for the completed basis. Those are proved in
different Lean files and have not been converted here into a
continuous conventional proof.

\section{The determinant mechanism}
\label{sec:determinant}
The intermediate result \texttt{one\_determinant} in
\leanfile{New/Heights/OneDeterminant.lean}%
{OneDeterminant.lean} explains the $2\delta$ in
\eqref{eq:ratio-two}. Begin with a basis $x_i$ of the row
space and the Gram determinant of its untwisted lifts.
The complex covolume assumption bounds every conjugate of
this determinant by $e^V$. At a prime $\ell$, a matrix
$M_\ell$ recombines the twisted lifts into integral vectors.
If $e_\ell$ is the local index loss and $\gamma_\ell$ a
scalar correction, the local determinant condition has the
shape
\begin{equation}
\label{eq:det-local}
|\det M_\ell|_\ell
 \ge e^{-e_\ell}|\gamma_\ell|_\ell^k.
\end{equation}
Integrality bounds the $\ell$-adic norm of the Gram
determinant after recombination by $1$. The determinant
changes by $\det M_\ell$ and its conjugate, so
\eqref{eq:det-local} costs at most $2e_\ell$ in the
logarithmic height. The sum of the $e_\ell$ over every
finite set is at most $\delta$. The scalar corrections
contribute their separate height budget. Combining the
complex and finite-place estimates gives the abstract
bound $V+2\delta+k g$ of \texttt{one\_determinant},
where $g$ bounds the scalar height.

\begin{proposition}[Abstract one-determinant bound]
\label{prop:one-determinant}
Let $W$ be a $D$-dimensional complex vector space with
basis $y_1,\ldots,y_D$. Let $u_0$ be an injective lift,
and let $u_\iota$ be a lift for each prime-adic
isomorphism $\iota$. Assume every complex conjugate of
$\det\operatorname{Gram}(u_0y)$ has modulus at most
$e^\alpha$. At each prime $\ell$, suppose there is a
real $e_\ell$ and, for every $\iota$ above $\ell$,
matrices recombining $u_\iota y$ into vectors with
integral pairings. Suppose their determinants have
norm at least $e^{-e_\ell}|\gamma_\iota|_\ell^k$
for a nonzero scalar family $\gamma_\iota$, and that
$\sum_{\ell\in U}e_\ell\le\delta$ for every finite
set of primes $U$. If
$\operatorname{HeightLe}(g,|\gamma|^{-2})$, then
\[
\operatorname{HeightLe}\!\left(\alpha+2\delta+kg,
\frac{\operatorname{vol}(u_\iota)^2}
     {\operatorname{vol}(u_0)^2}\right).
\]
\end{proposition}
\begin{proof}
At an $\ell$-adic embedding, pair each recombined
vector with the corresponding recombination obtained
after complex conjugation. The pairing hypothesis
makes the resulting Gram matrix integral, so the
$\ell$-adic norm of its determinant is at most one.
The determinant transformation law contributes one
factor from each recombination matrix. Their lower
bounds give a local upper bound
\[
\left|\det\operatorname{Gram}(u_\iota y)\right|_\ell
 \le e^{2e_\ell}|\gamma_\iota|_\ell^{-2k}.
\]
The untwisted Gram determinant is nonzero by
injectivity. Use it as part of the common clearing
factor in \texttt{HeightLe}. The complex estimate
contributes $\alpha$, the finite determinants contribute
$2e_\ell$, and the factor $|\gamma|^{-2k}$
contributes $kg$ by the power rule. The finite-set
bound gives total cost at most
$\alpha+2\delta+kg$.
\end{proof}
The formal declaration is \path{one_determinant} in
\leanfile{New/Heights/OneDeterminant.lean}%
{OneDeterminant.lean}. Its source type specifies the
integral lattices for both $\iota$ and its conjugate,
the matrix entries, and the exact pairing. The
proposition records the determinant calculation after
those local hypotheses have been established.

The independent rows from Theorems A and B need not be a
full basis. The result \texttt{Data.crude\_rows} in
\leanfile{New/Exchange/CrudeRows.lean}{CrudeRows.lean}
supplies additional forms that span the missing directions.
Their twisted lifts are integral after a controlled scalar
change, and all complex conjugates of their untwisted lifts
have squared norm at most $\exp(x_1^2b^2)$. The proof of
\texttt{heightLe\_volume\_rows} extends the selected rows
with these forms, applies the determinant estimate, and
absorbs the completion cost into $D_{\rm row}b$. The
last absorption uses the admissible scale inequalities;
it is a distinct quantitative step from the determinant
identity itself.

\section{Completing the basis with crude rows}
\label{sec:crude-rows}
The theorem \texttt{Data.crude\_rows} has more content
than mere spanning. It supplies a nonzero scalar family
$c(\vartheta)$ with a controlled square cost and a
spanning set of rows $y$ for which, whenever $\ell$ does
not divide the twist conductor,
\begin{equation}
\label{eq:crude-integral}
\frac{c(\vartheta)}{c(1)}
\,\iota_\vartheta y
\quad\text{is integral at every }\ell\text{-adic embedding}.
\end{equation}
Every conjugate of the untwisted lift has pointwise
squared norm at most $\exp(x_1^2b^2)$. Both properties
are in \leanfile{New/Exchange/CrudeRows.lean}%
{CrudeRows.lean}, declaration \texttt{Data.crude\_rows}.

The construction begins with a grid period for a form
at the far end of the theta chain. The declaration
\texttt{exists\_isTidy\_integral\_eq\_gridPeriod}
gives tidy finite data of order $4m$ and expresses a
pairing with the lifted form as a scalar multiple of
that period. The route theorem
\texttt{gridPeriod\_climb} carries the identity through
the $m$ stages. Its scalar factors have square cost
$2j^2$ through stage $j$; at $j=m$ this explains the
$2m^2$ cost in \texttt{crude\_rows}.
The separate declaration
\texttt{exists\_mul\_endPeriod\_mem\_embValRing}
clears the endpoint denominator by an integer no larger
than $\exp(x^{99}b^2)$. Finally a surjection from
grid-level forms and a dual-pairing argument show that
the resulting rows span. The source proves these period
and duality statements; the present description does not
replace their proofs.

There is a simple numerical reason that the coarse bound
is enough. Theorem B leaves at most
$4D_{\rm row}/b^2$ basis directions unfilled.
The logarithm of the squared norm of each crude row is
at most $x_1^2b^2$. Their total direct contribution is
therefore at most $4D_{\rm row}x_1^2$.
The hierarchy $x_1^{10^6}\le b$ makes this much smaller
than $D_{\rm row}b$. The actual source estimate also
absorbs scalar costs and determinant cross terms; its
final statement is \texttt{heightLe\_volume\_rows},
not the direct norm estimate alone.

\chapter{Dependencies and verification boundary}
\label{ch:boundary}
\section{A source map}
The table gives the reader a route into the pinned compact
source. Each row names a mathematical output and the Lean
declaration that supplies it. The table records source
locations; it does not certify the argument in those files.

\begin{longtable}{p{0.35\textwidth}p{0.57\textwidth}}
\toprule
Output & Lean declaration and file \\
\midrule
\endhead
Five target theorems & \texttt{Main.lean}, with independent
  statements in \texttt{Statements.lean} \\
Admissible data & \texttt{exists\_data},
  \texttt{New/Data/DataExist.lean} \\
Short forms and index & \texttt{Data.short\_forms},
  \texttt{New/Endgame/ShortAndEnough.lean} \\
Almost spanning translates & \texttt{Data.enough\_forms},
  \texttt{New/Endgame/ShortAndEnough.lean} \\
Family of bounded covolume &
  \texttt{Data.exists\_family\_covolumeSq\_le},
  \texttt{New/Endgame/LongForms.lean} \\
Lower covolume bound & \texttt{Data.integral\_forms\_long},
  \texttt{New/Endgame/LongForms.lean} \\
Weighted error identity & \texttt{Data.remainder\_eq},
  \texttt{New/Endgame/ShortAndEnough.lean} \\
Window point & \texttt{IsPolar.exists\_window\_sub\_mul\_le},
  \texttt{New/Endgame/OneSidedCriterion.lean} \\
Main inequality and GRH & \texttt{star\_inequality},
  \texttt{eventually\_primeError\_le\_rpow},
  \texttt{New.grh} in \texttt{New/Main.lean} \\
\bottomrule
\end{longtable}

The compact source separates \texttt{src/New}, the proof-specific
argument, from \texttt{src/Lib}, the background developed beyond
Mathlib. Its authors count 219 files under \texttt{New} and 573
under \texttt{Lib} \cite[Layout]{Sernau2026}. The guide names
class field theory, cohomology, Galois representations, $p$-adic
comparison, and automorphic formulas among the background inputs.
The distinction is organizational: the Lean proof includes these
results and does not import them as unproved axioms.

\section{What has been checked}
Sutherland's record pins the compact source to commit
\texttt{72fead10}. It reports that both the compact and readable
editions build, that their five statements match, and that the
theorems depend only on \texttt{propext},
\texttt{Classical.choice}, and \texttt{Quot.sound}. The
record includes replays by Lean's kernel and independent
checkers, with a successful paranoid Comparator run
\cite[Summary]{Sutherland2026}. These are checks of formal
derivations and their exact statements. They are distinct from
a mathematical referee reading of the argument or a validation
of this English exposition.

\begin{question}
\label{q:human-proof}
Can every substantive lemma on the compact dependency route,
including the background results in \texttt{src/Lib}, be matched
to a conventional proof with all hypotheses and normalizations
checked independently of the Lean development?
\end{question}
The compact authors explicitly say that this work has not yet
been done \cite[Status]{Sernau2026}. This manuscript gives a
route for that reading and leaves Question~\ref{q:human-proof}
open.

\appendix
\chapter{Source and AI methods}
\label{ch:methods}
This manuscript was drafted on 27 November 2026 by an OpenAI
Codex agent working for Taylor Dupuy. The child-agent interface
did not expose its model identifier or token telemetry. The
draft used the pinned compact source, its guide and README,
the earlier book's teaching architecture, and the verification
record cited above. It did not run Lean, inspect every one of
the compact source files, or independently referee the deep
mathematics. Its original contribution is expository text,
elementary deductions from the source interfaces, and the
source map. These require human review before they are
treated as a faithful account of the complete argument.

The run plan, exact selected sources, and source boundaries are
recorded in the working repository under
\texttt{ai/2026-11-27-grh-compact-book/}. A source file linked
from the body is a claim about what that file states, not a
claim that the English explanation proves its theorem.

\bibliographystyle{amsplain}
\bibliography{grh-compact}
\end{document}
% <<< [agent-codex 2026-11-27 compact-book]
